\section{总结与延伸：Generalist Agent 是一个闭环系统，不是单个模型}

整堂课从主动小猫出发，最终形成一套可迁移的系统分解：开放环境产生未穷尽的任务，互联网知识提供人类先验，foundation model 处理语言与多模态上下文，reward/code/action interface 把先验接到行为，critic 和 environment feedback 判断后果，skill memory 与 curriculum 支持长期积累。任何单点 benchmark 都只能验证其中一部分。

\subsection{课程收束：MineDojo、Agent Learning 与 Multimodal Prompts}

结尾三页回顾论文、agent demo 与 multimodal prompt，并把这些工作放回共同目标。它们不是相互替代的“产品列表”：MineDojo 提供 environment/data substrate，MineCLIP 学语义 reward，Voyager 学可执行 memory，Eureka 搜索 reward program，VIMA/RT-2/RoboCat 则探索统一多模态 action model。

\lecturefigure{slide-58-minedojo-paper-overview.jpg}{MineDojo 论文回顾：开放环境与互联网知识基础设施}{官方录像恢复幻灯片 V058；00:46:24}

\begin{knowledgebox}{读图：把 collage 还原为系统职责}
页面同时出现环境截图、互联网素材与论文信息。先辨认 environment、knowledge database 和 benchmark 三层，再问每层提供什么 evidence；它支持平台覆盖广，却不能单凭 collage 判断 policy quality。真正的验收来自任务协议、reward、baseline 与可复现 artifact。
\end{knowledgebox}

\lecturefigure{slide-59-agent-learning-examples.jpg}{开放世界 agent learning 的多种行为与任务示例}{官方录像恢复幻灯片 V059；00:46:27}

示例展示 agent 能在不同场景执行采集、制作、探索或交互。读图时不要把多张成功截图相加成 success rate；它们负责说明 task diversity，定量能力仍应看 held-out evaluation、failure distribution 与预算。开放世界尤其需要报告 agent 没有做成什么。

\lecturefigure{slide-60-multimodal-prompts-summary.jpg}{Multimodal prompt 统一视觉目标、语言关系与机器人任务}{官方录像恢复幻灯片 V060；00:46:51}

这页回到 VIMA 的核心接口：prompt 不是固定 instruction string，而是可组合的多模态 task program。先看 token 类型与对象 binding，再看模型输出 action；它支持 specification flexibility，却也增加 prompt ambiguity、adversarial input 和 safety validation 的责任。

\subsection{Grand Challenge：Human-level Minecraft AI}

课程用 human-level Minecraft 作为 community challenge 收尾。这里的“human-level”不能缩减为击败某个 boss 或获得一组物品，而应包含语言理解、长期记忆、创意建造、社交协作、失败恢复、跨版本适应与在新目标上持续学习。它是研究方向，不是课堂声称已经达到的里程碑。

\lecturefigure{slide-61-human-level-minecraft-challenge.jpg}{Grand Challenge：Human-level Minecraft AI}{官方录像恢复幻灯片 V061；00:47:18}

\begin{warningbox}{不要让“Human-level”被一个 leaderboard 偷换}
单项速度、地图覆盖或物品数量只能测局部能力。更完整的评价应分开报告 task proposal、knowledge use、planning horizon、memory retention、creative outcome、robustness、safety 和 human collaboration，并在未见 seed、版本和目标上测试。
\end{warningbox}

\teachervoice{讲者最后把 human-level Minecraft 留给整个社区，并半开玩笑说这也是“在工作中玩 Minecraft”的理由。真正的教学含义是：当前工作只触及开放世界 agent 的表面，下一阶段需要更强视频学习、多模态基础模型与主动闭环，而不是把一组漂亮 demo 当成终点。}

\subsection{八维系统检查表}

为了把本讲迁移到其他 agent 项目，可以对任何系统逐项审计以下八个维度。它们覆盖 problem、data、model、interface、learning loop 与 evidence，能够防止团队把 foundation model 的语言流畅度误当成完整的行动能力。

\begin{longtable}{@{}L{0.18\textwidth}L{0.34\textwidth}L{0.36\textwidth}@{}}
\toprule
维度 & 核心问题 & 典型失败 \\
\midrule
Environment & 世界是否有可组合规则与长期状态？ & 封闭 benchmark 被误称开放世界 \\
Goal & 目标如何提出、消歧并判断完成？ & 模糊 prompt 被 reward shortcut 利用 \\
Knowledge & 哪些先验来自 web、video、human 或 simulator？ & 来源噪声与当前状态不匹配 \\
Interface & 输出是 action、code、reward 还是 skill？ & 抽象层级与任务动力学不匹配 \\
Feedback & 谁提供可验证的 success、error 与 constraint？ & critic 自说自话、reward 未校准 \\
Memory & 成功经验怎样存储、检索、更新与失效？ & 错误 skill 长期污染后续任务 \\
Curriculum & 下一任务怎样平衡可行性、新颖性与风险？ & 太难停滞、太易重复、只追 reward \\
Evidence & 指标证明什么，又没有证明什么？ & simulator 局部优势被外推为通用智能 \\
\bottomrule
\end{longtable}

\begin{importantbox}{把本讲压缩成一句话}
Generalist agent 不是“一个更大的模型”，而是把\textbf{开放环境、互联网先验、可执行接口、真实反馈、持久记忆与自动课程}连接成闭环；模型的价值取决于它能否在这个闭环里把 passive knowledge 变成经过验证的 active competence。
\end{importantbox}

\subsection{最值得保留的四个工程判断}

本讲虽然以 Minecraft 和机器人为例，但四个判断可以直接迁移到 tool-using agents、browser agents、coding agents 和 scientific agents。它们共同强调：系统能力由 environment/interface/evidence 决定，而不是由模型名称单独决定。

\begin{enumerate}
  \item \textbf{先选 action abstraction。} 高频 motor control、API call、程序和 reward specification 需要不同接口；不要让模型承担不适合其时间尺度的职责。
  \item \textbf{把 execution result 当成 truth source。} 语言自评只能辅助，真正的 success 应来自环境状态、测试、constraint 与人工审查。
  \item \textbf{把 memory 存成可验证 artifact。} 代码、skill、trajectory 和 reward program 应带前置条件、版本、测试与失效策略。
  \item \textbf{把指标限制写进结论。} 曲线和 demo 只支持其评估协议内的主张；跨环境、跨 embodiment 与安全能力需要独立证据。
\end{enumerate}

\subsection{拓展阅读}

\begingroup
\footnotesize
\begin{itemize}[itemsep=0pt,topsep=0.15em,parsep=0pt]
  \item Fan et al., \href{https://arxiv.org/abs/2206.08853}{\emph{MineDojo: Building Open-Ended Embodied Agents with Internet-Scale Knowledge}}
  \item Wang et al., \href{https://arxiv.org/abs/2305.16291}{\emph{Voyager: An Open-Ended Embodied Agent with Large Language Models}}
  \item Ma et al., \href{https://arxiv.org/abs/2310.12931}{\emph{Eureka: Human-Level Reward Design via Coding Large Language Models}}
  \item Jiang et al., \href{https://arxiv.org/abs/2210.03094}{\emph{VIMA: General Robot Manipulation with Multimodal Prompts}}
  \item Baker et al., \href{https://arxiv.org/abs/2206.11795}{\emph{Video PreTraining (VPT): Learning to Act by Watching Unlabeled Online Videos}}
  \item Brohan et al., \href{https://arxiv.org/abs/2307.15818}{\emph{RT-2: Vision-Language-Action Models Transfer Web Knowledge to Robotic Control}}
  \item Reed et al., \href{https://arxiv.org/abs/2306.11706}{\emph{RoboCat: A Self-Improving Foundation Agent for Robotic Manipulation}}
\end{itemize}

阅读这些工作时，可以反复应用八维检查表，并额外追问：训练数据是否包含 action consequence，evaluation 是否覆盖 failure，memory 是否会累积错误，reward 是否可能被钻漏洞，接口是否适合当前 embodiment，以及结论是否越过了 simulator、任务集合与时间边界。
\endgroup

\end{document}
